\documentclass[11pt]{article}

\usepackage[margin=1in]{geometry}
\usepackage{amsmath, amssymb, amsthm}
\usepackage{bbm}
\usepackage{enumitem}

\title{Mathematics 629 \\ Spring 2026 \\ Homework Assignment 5}
\date{Due Wednesday, March 11}
\author{}

\begin{document}
\maketitle
\thispagestyle{empty}

\begin{enumerate}[leftmargin=*, label=\textbf{\arabic*.}]

	%% Problem 1
	\item Let $f, g : X \to \mathbb{R}$ be integrable functions. Assume that
	      \[
		      \int_E f\, d\mu \leq \int_E g\, d\mu
	      \]
	      for all measurable sets $E$. Show that $f \leq g$ a.e.

	      \begin{proof}
		      First let $h(x) = f(x) - g(x)$. Then by theorem 1.115
		      $$\int_E f d\mu - \int_E g d\mu = \int_E h d\mu \leq 0.$$
		      Let $E = \{x: h(x) > 0\}$. Equivalently, $E = \{x: f(x) > g(x) \}$ which is measurably by corollary 1.75.
		      Notice that by corollary 1.83, $h^+ = max\{h,0\}$ is measurable and by the construction of $E$ that
		      $$\int_E h d\mu = \int_E h^+ d\mu \leq 0.$$
		      Since $h^+ \geq 0$ then $\int_E h^+ d\mu = 0$ so by theorem 1.104 $h^+(x) = 0$ a.e. In other words, $h(x) \leq 0$ almost everywhere so $f(x) \leq g(x)$ almost everywhere.

	      \end{proof}

	      %% Problem 2
	\item We use Egorov's theorem to prove the dominated convergence theorem, but it can also be deduced by some trickery with Fatou's lemma.

	      Suppose $f_n, g_n, f, g$ are $L^1$ functions such that $\int g_n\, d\mu \to \int g\, d\mu$, $|f_n| \leq g_n$, $f_n \to f$ a.e., and $g_n \to g$ a.e.

	      \begin{enumerate}[label=(\roman*)]
		      \item Show that $\displaystyle\lim_{n\to\infty} \int f_n\, d\mu = \int f\, d\mu$.

		            \textit{Hint:} Assume that all functions are real-valued and that the a.e.\ convergence assumptions are replaced by pointwise convergence assumptions. Apply Fatou's lemma to the nonnegative sequences $g_n + f_n$, $g_n - f_n$ and derive the inequalities
		            \[
			            \int (g + f)\, d\mu \leq \int g\, d\mu + \liminf_{n\to\infty} \int f_n\, d\mu
		            \]
		            \[
			            \int (g - f)\, d\mu \leq \int g\, d\mu - \limsup_{n\to\infty} \int f_n\, d\mu
		            \]

		      \item With the same assumptions on $f_n, g_n, f, g$ as in (i), show that
		            \[
			            \int |f_n - f|\, d\mu \to 0.
		            \]

		      \item Deduce the dominated convergence theorem.

		      \item Show that if $f_n, f \in L^1$ and $f_n \to f$ a.e.\ then
		            \[
			            \int |f_n - f|\, d\mu \to 0 \iff \int |f_n|\, d\mu \to \int |f|\, d\mu.
		            \]
	      \end{enumerate}


	      %% Problem 3
	\item Let $I = [0, \infty)$, $m$ Lebesgue measure on $I$. Let
	      \[
		      f(x) = e^{x/2} \cos(e^x).
	      \]
	      Show that $\displaystyle\lim_{R\to\infty} \int_0^R f(x)\, dm$ exists but $f \notin L^1(I, m)$.
	      \begin{proof}
		      Let $f(x) = e^{x/2}cos(e^x)$. By the note for problems 3 and 4,
		      $$ \lim_{R\to \infty} \int_0^R f(x) dm = \int_{[0,\infty]}f(x) dm = \int_0^\infty f(x)dx.$$
		      To prove $f \notin L^1(I,m)$ we will show $\int_0^\infty \vert f\vert dm = \infty.$ Let $u = e^x$ and $du = e^x dx$. By u-sub we get
		      $$ \int_1^\infty u^{-1/2}\vert cos(u)\vert du.$$ Observe that $[0,\infty] \supset \uplus_{k=1}^\infty [k\pi - \frac{\pi}{4}, \pi + \frac{\pi}{4}]$ so by monotonicity,
		      $$\int_1^\infty u^{-1/2}\vert cos(u) \vert du \geq \sum_{k=1}^\infty \int_{k\pi - \frac{\pi}{4}}^{k\pi+\frac{\pi}{4}} u^{-1/2}\vert cos(u) \vert du.$$
		      Let $\vert t \vert \leq \frac{\pi}{4}$. Since $\vert cos(k\pi + t)\vert = \vert \cos(t)\vert \geq cos(\frac{\pi}{4}) = \frac{1}{\sqrt{2}}$ and $u^{-1/2} \geq (k\pi + \frac{\pi}{4}$ then
		      $$ \int_{k\pi - \frac{\pi}{4}}^{k\pi+\frac{\pi}{4}} u^{-1/2}\vert cos(u) \vert du \geq \frac{1}{\sqrt{2}}(k\pi + \frac{\pi}{4})^{-1/2} \frac{\pi}{2}.$$ By bounding $(k\pi + \frac{\pi}{4})^{-1/2} \geq (2k\pi)^{-1/}$ then
		      $$\sum_{k=1}^\infty \int_{k\pi - \frac{\pi}{4}}^{k\pi+\frac{\pi}{4}} u^{-1/2}\vert cos(u) \vert du \geq \sum_{k=1}^\infty \frac{\sqrt{\pi}}{4\sqrt{k}}$$
		      which diverges by the p-series test, thus we have shown $\int_0^\infty \vert f\vert dm = \infty.$


	      \end{proof}

	      %% Problem 4
	\item Let $f \in L^1(\mathbb{R}, m)$. Prove that
	      \[
		      \lim_{\lambda \to \infty} \int_{\mathbb{R}} \cos(\lambda x)\, f(x)\, dm = 0.
	      \]
	      \begin{proof}
		      Let $\epsilon >0$ and choose a compactly supported $g$ such that
		      $\verts f - g\verts_1 < \epsilon/2$ and for $R >0$, $supp(g) \subset (-R, R)$. Since $|cos(\lambda x)| \leq 1$ then by theorem 1.117
		      $$| \int_\mathbb{R} cos(\lambda x)f(x)dm - \int_\mathbb{R} cos(\lambda x)g(x) dm | \leq \int_\mathbb{R} |f-g|dm < \epsilon/2.$$
		      Since $supp(g) \subset (-R,R)$ then
		      $$\int_\mathbb{R} cos(\lambda x)g(x)dm = \int_{-R}^R cos(\lambda x)g(x) dx$$ by the hint. Integrating by parts shows
		      $$\int_{-R}^R cos(\lambda x)g(x)dx = \frac{sin(\lambda x)}{\lambda}g(x) \vert_{-R}^R - \int_{-R}^R \frac{sin(\lambda x)}{\lambda}g'(x)dx.$$ By construction of the compact support of $g$,
		      $$\int_{-R}^R cos(\lambda x)g(x)dx \leq |\int_{-R}^R \frac{sin(\lambda x)}{\lambda}g'(x)dx | \leq \frac{c}{\lambda} < \frac{\epsilon}{2}$$
		      for some $c \in \mathbb{R}$ as $\lambda \to \infty$. Thus
		      $$|\int_\mathbb{R} cos(\lambda x)f(x)dm| \leq |\int cos(\lambda x)(f-g)dm| + |\int cos(\lambda x)gdm| < \frac{\epsilon}{2} + \frac{\epsilon}{2} = \epsilon$$
		      so $\lim_{\lambda \to \infty} \int_{\mathbb{R}} \cos(\lambda x)\, f(x)\, dm = 0$ as a result.


	      \end{proof}

	      \medskip
	      \noindent\textit{Note for problems 3 and 4:} Here you are allowed to use that if $g$ is Riemann integrable on a bounded interval $[a,b]$ then the Riemann integral $\int_a^b f(x)\, dx$ coincides with the Lebesgue integral $\int_{[a,b]} f\, dm$, a fact that we will establish in class.

	      %% Problem 5
	\item Determine the limits
	      \begin{enumerate}[label=(\roman*)]
		      \item $\displaystyle\lim_{n\to\infty} \int_0^{n^{99/100}} \left(\frac{x}{n}\right)^n dx$

		      \item $\displaystyle\lim_{n\to\infty} \int_0^n \left(1 + \frac{x}{n}\right)^n e^{-2x}\, dx$
	      \end{enumerate}
	      and, in both cases, carefully justify your computation.

	      \begin{enumerate}
		      \item Notice when $x = 0, (\frac{x}{n})^n = 0$ and that for all $n$, plugging in $n^{\frac{99}{100}}$ gives
		            $\frac{n^{\frac{99}{100}}}{n} = n^{\frac{-1}{100}}$. In other words, $(\frac{x}{n})^n$ is nonnegative, bounded and monotonically decreasing as $n\to \infty$. By theorem 1.121, it is sufficient to directly compute the integral which gives
		            $$ \lim_{n\to\infty}\int_0^{n^{99/100}} (\frac{x}{n})^n dx = \lim_{n\to\infty}\int_0^{n^{-1/100}}u^n ndu = \lim_{n\to \infty} \frac{n}{n+1}n^{-\frac{n+1}{100}} = 0$$
		            via u-substitution ($u = \frac{x}{n}$, $dx = ndu$) and the power rule.
		      \item To use DCT, $\{f_n\}$ must be a bounded sequence of non-negative functions that converges to $f$. Let $f_x = (1 + \frac{x}{n})^n e^{-2x}.$ Observe that
		            $(1 + \frac{x}{n})^n \leq e^x$ since $n\ln(1+\frac{x}{n}) \leq n \frac{x}{n} = x = \ln(e^x).$ Since $e^{-2x}  \leq 1$ on $[0,n]$ then $f_n$ is bounded.

		            Then use the taylor expansion of $n\ln(1+\frac{x}{n}) = x + \frac{x^2}{2n} + \frac{x^3}{3n^2} + \cdots$ to see
		            $\lim_{n\to \infty} n\ln(1+\frac{x}{n}) = x.$ Thus $e^{\ln(1+\frac{x}{n})} \to e^x$ as $n\to \infty$ so by substitution,
		            $(1 + \frac{x}{n})^n e^{-2x} \to e^xe^{-2x} = e^{-x}.$ By theorem 1.127,
		            $$\lim_{n\to \infty} \int_0^n \left(1 + \frac{x}{n}\right)^n e^{-2x}\, dx = \int_0^\infty e^{-x}\, dx = 1.$$


	      \end{enumerate}

\end{enumerate}
\end{document}
